\documentclass[11pt,a4paper]{article}
\usepackage[T1]{fontenc}
\usepackage[utf8]{inputenc}
\DeclareUnicodeCharacter{2124}{\ensuremath{\mathbb{Z}}}
\DeclareUnicodeCharacter{0394}{\ensuremath{\Delta}}
\DeclareUnicodeCharacter{2265}{\ensuremath{\geq}}
\DeclareUnicodeCharacter{00B2}{\ensuremath{{}^2}}
\usepackage{lmodern}
\usepackage{amsmath,amssymb,amsthm,mathtools}
\usepackage{geometry}
\geometry{margin=2.6cm}
\usepackage{booktabs,array}
\usepackage{enumitem}
\usepackage[dvipsnames]{xcolor}
\usepackage[hidelinks]{hyperref}
\usepackage{xurl}
\hypersetup{pdftitle={The gap between the two walls is positive},pdfauthor={Toshihisa Urahigashi}}
\setlist{itemsep=2pt,topsep=4pt}
\raggedbottom
\emergencystretch=3em

\newcommand{\Z}{\mathbb{Z}}
\newcommand{\Q}{\mathbb{Q}}
\newcommand{\Zp}{\mathbb{Z}_p}
\newcommand{\Qp}{\mathbb{Q}_p}
\newcommand{\I}{\mathcal{I}}
\newcommand{\Ocal}{\mathcal{O}}
\newcommand{\ordp}{\operatorname{ord}_p}
\newcommand{\fr}{\operatorname{frac}}
\newcommand{\hull}{\operatorname{hull}}
\newcommand{\mulog}{\mu^{\log}}
\newcommand{\GL}{\mathrm{GL}}
\newcommand{\Vbad}{\mathbb{V}_{\mathrm{mod}}^{\mathrm{bad}}}
\newcommand{\qb}{\underline{q}}
\newcommand{\itwo}{\mathrm{ind2}}
\newcommand{\leanus}{\textunderscore\allowbreak}
\newcommand{\lean}[1]{{\ttfamily\let\_\leanus #1}}

% labels
\newcommand{\lab}[2]{{\upshape\fcolorbox{#1}{#1!8}{\scriptsize\sffamily\bfseries\textcolor{#1}{#2}}}}
\newcommand{\tagS}{\lab{blue}{S}}
\newcommand{\tagX}{\lab{Purple}{X}}
\newcommand{\tagP}{\lab{ForestGreen}{P}}
\newcommand{\tagR}{\lab{teal}{R}}
\newcommand{\tagN}{\lab{BrickRed}{N}}
\newcommand{\tagH}{\lab{BurntOrange}{H}}
\newcommand{\known}[1]{\par\noindent{\footnotesize\upshape Known mathematics: #1}\par\smallskip}
\newcommand{\q}[1]{``#1''}
\newcommand{\xquote}[2]{\begin{quote}\small\tagX\ #1\par\nopagebreak\hfill{\footnotesize #2}\end{quote}}

\theoremstyle{plain}
\newtheorem{theorem}{Theorem}[section]
\newtheorem{proposition}[theorem]{Proposition}
\newtheorem{corollary}[theorem]{Corollary}
\newtheorem{lemma}[theorem]{Lemma}
\theoremstyle{definition}
\newtheorem{definition}[theorem]{Definition}
\newtheorem{example}[theorem]{Example}
\theoremstyle{remark}
\newtheorem{remark}[theorem]{Remark}

\title{The gap between the two walls is positive:\\ a closed form for the (Ind2)-terms of lana-agents/iut}
\author{Toshihisa Urahigashi\\\small ORCID 0009-0004-0460-6242\\\small\texttt{123026.urahigashi.toshihisa@gmail.com}}
\date{Version 1.1.0, 8 November 2026\\\small\href{https://doi.org/10.5281/zenodo.24039377}{doi:10.5281/zenodo.24039377}}

\begin{document}
\maketitle

\begin{abstract}
The public Lean formalization lana-agents/iut writes the comparison step in the proof of Corollary~3.12 of
inter-universal Teichm\"uller theory (IUT)~III---the ``wall''---in two forms: a variant of Corollary~3.12, (Var),
and the Scholze--Stix inequality, (SS). The difference $\Delta(D)$ of their right-hand sides is a weighted average of
(Ind2)-terms at the bad residue characteristics. The repository proves $\Delta(D)\ge0$ and records as open whether
these terms are zero or positive.

We determine these terms in closed form. The (Ind2)-term equals $(\ordp\theta-\kappa_v(\theta))\log p$, where
$\theta=\qb_v^{\,j^2}$ is the theta value and $\kappa_v(\theta)$ is the largest integer $k$ with
$\theta\,\I_v\subseteq p^k\I_v$, $\I_v$ being the log-shell. Consequently the term is always at least
$\fr(\ordp\theta)\log p$, with equality when $p$ is odd and $e_v\le p-2$ (there the log-shell is a fractional
ideal). At a bad place $\ordp\theta$ is never an integer, because $\ell$ does not divide the order of $q$.
Hence $\Delta(D)>0$ for every initial $\Theta$-datum, and in this formalization (SS) is strictly stronger than (Var).

On the other hand $\Delta(D)$ is bounded above by a quantity that does not involve the size of the $q$-orders (the
values $\log p$ at the bad places and the ramification indices). The two walls are height inequalities of the same
shape that differ by a term of conductor type. The average over the labels is given by the class number of
$\Q(\sqrt{-\ell})$. We compute the values for six explicit initial $\Theta$-data.

No claim is made about the truth of (Var), (SS), Corollary~3.12, or the abc conjecture.
\end{abstract}

\tableofcontents

\section{Introduction}

\subsection{Motivation}
The ultimate aim of this work is to understand exactly the algorithm by which IUT computes the height inequality.
The computation of IUT~IV combines, step by step, theta values, log-shells, the indeterminacies (Ind1), (Ind2),
(Ind3), holomorphic hulls and log-volumes, and produces the two sides of the height inequality. To write this
algorithm down exactly, including the value produced at each step, every step must have a determined value. The
(Ind2)-term, whose value the public formalization lana-agents/iut leaves open, is one of these steps. This note
determines it.

\subsection{The two walls and their gap}
The repository lana-agents/iut~\cite{lana} (``lana'' below) is a public Lean formalization that builds the objects of
IUT. Its directory \lean{Iut/WallEquiv/} collects formulations of the step from Theorem~3.11 to Corollary~3.12 of
IUT~III~\cite{IUT3}, the \q{wall} in the words of the interim report of Project LANA~\cite[\S0.2]{LANA}, and proves
equivalences between them in its models. The file \lean{Gap.lean} treats the difference between two of them.

\xquote{``Christian's question: \emph{is the LANA wall any weaker than the Scholze--Stix wall?}''\quad
``\textbf{Open}: whether the (Ind2)-terms at the bad primes are \lean{0} or \lean{> 0} (the indeterminacy
automorphisms at a prime are all \lean{ℤ\_p}-lattice automorphisms of the factor log-shells, \lean{tensorAut}, which
need not commute with the theta values); so the sign \lean{Δ(D) ≥ 0} is all that is proved.''}{\cite{lana},
\lean{Iut/WallEquiv/Gap.lean}, module docstring}

For initial $\Theta$-data $D$ the two walls are the following inequalities (notation in Section~\ref{s:setting}).
\begin{itemize}
\item[(Var)] $-|\log q|\le R$, where $R=-|\log\Theta|$ is the log-volume of the holomorphic hull of the
$\Theta$-pilot region: a variant of Corollary~3.12.
\item[(SS)] $-|\log q|\le\mulog(\mathrm{LGP}\cdot\hull(VC(I)))$: the Scholze--Stix inequality~\cite{SS}, written with
IUT's container $\hull(VC(I))$.
\end{itemize}
lana defines the gap $\Delta(D)=R-\mulog(\mathrm{LGP}\cdot\hull(VC(I)))\ge0$, proves that (SS) is equivalent to
$-|\log q|\le R-\Delta(D)$ (\lean{ss\_iff\_le\_rhs\_sub\_gap}), and proves that $\Delta(D)$ is a weighted average of
the (Ind2)-terms at the bad residue characteristics (\lean{gap\_eq}). The same repository expects positivity:

\xquote{``they do not commute with the multiplication by \lean{q\^{}\{j²\}} in general, so \lean{ind2Term > 0} is to be
expected when \lean{ord(q\^{}\{j²\})} is not a multiple of the ramification index, but this is not
formalized.''}{\cite{lana}, \lean{Iut/WallEquiv/README.md}, section ``The gap between (Var) and (SS) (\lean{Gap.lean})''}

\subsection{Results}
Working with lana's definitions as they stand, we determine the (Ind2)-terms.
\begin{itemize}
\item Theorem~\ref{t:closed} (closed form): $\itwo=(\ordp\theta-\kappa_v(\theta))\log p$.
\item Theorem~\ref{t:bounds} (bounds): $\itwo\ge\fr(\ordp\theta)\log p$ and $\itwo<(1+a_v+b_v)\log p$, with equality in
the lower bound when $p$ is odd and $e_v\le p-2$; here $a_v,b_v$ are the constants of Proposition~1.2 of
IUT~IV~\cite{IUT4}.
\item Theorem~\ref{t:pos} (positivity): $\itwo>0$ at every bad place. Hence $\Delta(D)>0$ for every initial
$\Theta$-datum, and (SS) is strictly stronger than (Var).
\item Corollary~\ref{c:size} (size): $\Delta(D)<\sum_uw_u(\log p_u+3+\log e_u)$; the right-hand side does not involve
the size of the $q$-orders.
\item Proposition~\ref{p:classno} (label average): if $\ordp(\qb_v)=n/\ell$, the average of $\fr(\ordp\theta)$ over the
labels is $\tfrac12-\chi_\ell(n)\,h(-\ell)/(\ell-1)$ when $\ell\equiv3\bmod4$; at small places this is the average of
the term in units of $\log p$.
\end{itemize}
Our answer to lana's question is therefore: in lana's formalization the LANA wall is strictly weaker than the
Scholze--Stix wall, but only by a quantity of conductor type, which does not grow with the size of $q$.

\tagX\ The title of this note is close to the message of lana's commit \lean{6ef57cb} (13 October 2026),
\q{WallEquiv/Gap: the gap between (Var) and (SS) is positive for every D}. That commit proved the positivity of the
gap to the form of (SS) without the hull, from the archimedean defect (\lean{gap\_pos}); since the switch to the
hull in commit \lean{cdc9602} (14 October 2026) this theorem is \lean{gapNoHull\_pos}. With the hull on both sides
these defects cancel (\q{The shell and archimedean defects are exactly what the hull removes}, \lean{Gap.lean}), and
for $\Delta(D)$ the file proves the sign $\Delta(D)\ge0$ and leaves positivity open, as quoted above. The positivity in
the title of this note concerns $\Delta(D)$.

\subsection{Relation to earlier work}
\begin{itemize}
\item \tagX\ Positivity itself is expected, under a condition, in lana's documentation quoted above. We give a proof
and the exact value.
\item \tagX\ Dupuy and Hilado describe (Ind2) as acting through the $\Qp$-linear automorphisms arising from
$\Zp$-lattice automorphisms of $\I_v$~\cite[\S4.8]{DH1}, and bound the hull under such automorphisms \emph{from
above}~\cite[\S6.2]{DH2}; the floor $\lfloor\ordp(a)-\ordp(\beta)\rfloor$ appears in their estimate (6.6)--(6.7).
Theorem~\ref{t:closed} shows that the hull actually reaches the position of this rounding. As an alternative that
they do not pursue, they suggest writing out an explicit $\Zp$-basis and computing the action of these automorphisms
explicitly~\cite[Rem.~6.2.1(1)]{DH2}; Lemmas~\ref{l:orbit}--\ref{l:tensor} compute, for the lattices
$\theta\cdot\bigotimes\I$ of lana's definitions, the span of their images under these automorphisms.
\item \tagS\ Proposition~1.2(ii) of IUT~IV~\cite{IUT4} likewise bounds $\phi(p^\lambda(R_I)^{\sim})$ from above by
inclusions involving a floor, and (iii) replaces the floor by an allowance of $1$. Neither states an equality.
\item \tagP\ The lattice lemma (Lemma~\ref{l:orbit}) is elementary and not new. What is new, to the author's
knowledge, is its application to lana's definitions to determine the (Ind2)-terms, and the positivity, bounds and
label average that follow.
\end{itemize}

\subsection{What is not claimed}
\begin{itemize}
\item The object of this note is the definitions of lana's formalization. lana states that its (Ind2) group is a
\q{deliberate enlargement} that contains the image of $\mathrm{Ism}$ (Section~\ref{s:setting}). For the group of the
original papers the value lies between $0$ and the value found here; under the identification of
\cite[Prop.~9]{Fourth} it is $0$ (Section~\ref{s:groups}).
\item Nothing is claimed about the truth of (Var), (SS), Corollary~3.12 or the abc conjecture.
\item The proofs are on paper. They have not been checked in Lean on top of lana's definitions.
\end{itemize}

\subsection{Labels}
\tagS\ a statement of the original papers of IUT or of Mochizuki's documents. \tagX\ a statement of others (lana,
Project LANA, Dupuy--Hilado, Zeng). \tagP\ a statement proved here; \tagP\ alone marks an elementary fact. \tagR\ known
mathematics, applied here in a new way; the known mathematics is named. \tagN\ a result new in this note, to the
author's knowledge. \tagH\ a reading or interpretation, not a proof.

\section{Setting}\label{s:setting}
We follow lana's definitions; Lean names are given in parentheses.
\begin{itemize}
\item \emph{Initial $\Theta$-data} $D$: a number field $F$, an elliptic curve $E/F$, a prime $\ell\ge5$,
$\ell^*=(\ell-1)/2$, $K=F(E[\ell])$, the field of moduli $F_{\mathrm{mod}}$ and a nonempty set $\Vbad$ of bad places
(\lean{bad\_nonempty}). At a bad place $w$ the order of the Tate parameter $q_w$ is prime to $\ell$
(\lean{q\_order\_coprime}; IUT~I, Definition~3.1(c)~\cite{IUT1}).
\item \emph{Capsules and labels}: the labels of capsule $i=0,\dots,\ell^*-1$ are $\{0,\dots,i+1\}$, and the
distinguished label is $j=i+1$ (\lean{distinguished}). Thus $j$ runs through $1,\dots,\ell^*$.
\item \emph{Packets}: for a prime $p$ and a map (tuple) $c$ from the labels to the places of $F_{\mathrm{mod}}$ over $p$,
$T_c=\bigotimes_lK_{v_l}$, the tensor product over $\Qp$, where $v_l$ is the place of $K$ over $c(l)$ given by the
section of places.
\item \emph{Log-shells}: $\I_v=(2p)^{-1}\log_p(\Ocal_v^{\times})\subseteq K_v$, a $\Zp$-lattice
(\lean{factorLogShell}); $\bigotimes\I$ is the $\Zp$-submodule spanned by the pure tensors of log-shell elements
(\lean{LocalTheory.shell} $=$ \lean{tensorLogShell}).
\item \emph{Integral structure and hull}: $\Ocal$ is the integral closure of $\Zp$ in $T_c$ (\lean{integral}). The
hull $\hull(U)$ of a subset $U$ is the intersection of all $b\,\Ocal$, $b$ a unit, that contain $U$
(\lean{leastHull}; IUT~III, Remark~3.9.5(i)~\cite{IUT3}).
\item \emph{Log-volume}: $\mulog(S)=\log\bigl(\mu(S)/\mu(\Ocal)\bigr)/[T_c:\Qp]$ (\lean{componentVol}). Multiplication by
$p$ shifts it by $-\log p$, and multiplication by $y$ in the factor of label $l$ by $-\ordp(y)\log p$
(\lean{componentVol\_prod\_incl\_smul\_integral}), with $\ordp$ normalized by $\ordp(p)=1$.
\item \emph{The (Ind2) group}: $\mathrm{tensorAut}$ consists of the tensor products $\bigotimes_lg_l$ of
$\Qp$-linear automorphisms $g_l$ of the factors with $g_l(\I_{v_l})=\I_{v_l}$ (\lean{tensorAut}).
\item \emph{Theta values}: $\qb_v$ is a $2\ell$-th root of the Tate parameter at a bad place $v$, and $1$ elsewhere
(\lean{qrootOf}); $\theta=\qb_{v}^{\,j^2}$, with $v=v_j$, sits in the factor of label $j$ (\lean{scaleElt}). One has
$\ordp(\qb_v)=\operatorname{ord}_w(q_w)/(2\ell\,e_w)$, where $w$ is the place of $F$ below $v$ and $e_w$ the
ramification index of $F_w/\Qp$ (\lean{ordp\_qrootOf}).
\item \emph{The $\Theta$-pilot component} (in IUT~IV's reading):
$\Theta_{i,p,c}=\bigcup_{\phi\in\mathrm{tensorAut}}\phi\bigl(\theta\cdot\bigotimes\I\bigr)$
(\lean{thetaFinite\_eq}).
\item \emph{The (Ind2)-term} (\lean{ind2Term}):
\[\itwo(i,p,c)=\mulog\bigl(\hull(\Theta_{i,p,c})\bigr)-\mulog\bigl(\theta\cdot\hull(\textstyle\bigotimes\I)\bigr).\]
\item \emph{The gap} (\lean{gap\_eq}):
\[\Delta(D)=\frac1{\ell^*}\sum_i\sum_p\sum_c w_c\,\itwo(i,p,c),\]
with $w_c=\prod_lw(c(l))$ and $w(u)=[F_{\mathrm{mod},u}:\Qp]/[F_{\mathrm{mod}}:\Q]$.
The weights of the places over $p$ sum to $1$ (\lean{weight\_sum\_one}).
\end{itemize}

\xquote{``This class contains the image of \lean{Ism} and is in general larger: a \textbf{deliberate enlargement},
which makes the theta region larger and the hypothesis \lean{Cor312VariantHolds} weaker than with \lean{Ism}
exactly.''}{\cite{lana}, \lean{Iut/Concrete/LocalConstruct/IndAut.lean}, module docstring}

We use the constants of IUT~IV, Proposition~1.2~\cite{IUT4}, for the ramification index $e$ of $K_v/\Qp$:
\[a=\tfrac1e\lceil\tfrac{e}{p-2}\rceil\ (p>2),\quad a=2\ (p=2),\qquad
b=\Bigl\lfloor\frac{\log(pe/(p-1))}{\log p}\Bigr\rfloor-\frac1e .\]
By Proposition~1.2(i) of IUT~IV, $p^a\Ocal_v\subseteq\log_p(\Ocal_v^\times)\subseteq p^{-b}\Ocal_v$, with equality
in both inclusions when $p>2$ and $e\le p-2$; then $a=1/e=-b$.

\begin{definition}
A finite place $v\mid p$ of $K$ is \emph{small} if $p$ is odd and $e_v\le p-2$. At a small place
$\log_p(\Ocal_v^\times)=p^{1/e_v}\Ocal_v$, so $\I_v$ is an $\Ocal_v$-module.
\end{definition}

\section{Main results}

\begin{definition}
For a finite place $v\mid p$ of $K$ and $\theta\in K_v^{\times}$ put
\[\kappa_v(\theta)=\max\{k\in\Z:\ \theta\,\I_v\subseteq p^k\,\I_v\}.\]
The maximum exists because $\theta\,\I_v$ is again a $\Zp$-lattice.
\end{definition}

\begin{theorem}[closed form]\label{t:closed}\tagP\tagN\
Let $v=v_j$ and $\theta=\qb_v^{\,j^2}$. Then
\[\itwo(i,p,c)=\bigl(\ordp\theta-\kappa_v(\theta)\bigr)\log p .\]
In particular the term depends only on $j$ and $v_j$, not on the places at the other labels.
\end{theorem}

\begin{theorem}[bounds]\label{t:bounds}\tagP\tagN\
With the same notation and $e_v$ the ramification index of $K_v/\Qp$,
\[\fr(\ordp\theta)\,\log p\ \le\ \itwo(i,p,c)\ \le\ \bigl(\ordp\theta-\lfloor\ordp\theta-a_v-b_v\rfloor\bigr)\log p\
<\ (1+a_v+b_v)\log p .\]
If $v$ is small then $a_v+b_v=0$, the two bounds coincide, and $\itwo(i,p,c)=\fr(\ordp\theta)\log p$.
\end{theorem}

\begin{theorem}[positivity]\label{t:pos}\tagP\tagN\
\begin{enumerate}[label=(\roman*)]
\item If $v_j$ is a bad place, then $\ordp\theta\notin\Z$ for every $j=1,\dots,\ell^*$, and $\itwo(i,p,c)>0$.
\item $\Delta(D)>0$ for every initial $\Theta$-datum $D$. Hence in lana's formalization (SS) is strictly stronger than
(Var): (Var) holds and (SS) fails exactly when $-|\log q|\in(R-\Delta(D),R]$
(\lean{variant\_and\_not\_ss\_iff}), and this interval has positive length $\Delta(D)$.
\end{enumerate}
\end{theorem}

\begin{corollary}[formula and size of the gap]\label{c:size}\tagP\
For $u\in\Vbad$ let $v(u)$ be the place of $K$ over $u$, $r_u=\ordp(\qb_{v(u)})$, and $p_u$ the residue
characteristic. Then
\[\Delta(D)=\frac1{\ell^*}\sum_{j=1}^{\ell^*}\sum_{u\in\Vbad}w_u\,\bigl(j^2r_u-\kappa_{v(u)}(\qb_{v(u)}^{\,j^2})\bigr)\log p_u,\]
\[\frac1{\ell^*}\sum_{j=1}^{\ell^*}\sum_{u\in\Vbad}w_u\,\fr(j^2r_u)\log p_u\ \le\ \Delta(D)\ <\
\sum_{u\in\Vbad}w_u\bigl(\log p_u+3+\log e_{v(u)}\bigr).\]
If every bad place is small, the lower bound is an equality. The upper bound does not involve the size of the
$q$-orders.
\end{corollary}

\begin{proposition}[label average and class number]\label{p:classno}\tagP\tagR\
\known{Dirichlet's class number formula.}
Suppose $r_u=n_u/\ell$ with $n_u\in\Z$ and $\ell\nmid n_u$. Then
\[\frac1{\ell^*}\sum_{j=1}^{\ell^*}\fr\Bigl(\frac{j^2n_u}{\ell}\Bigr)=
\begin{cases}\dfrac12-\chi_\ell(n_u)\,\dfrac{h(-\ell)}{\ell-1}&(\ell\equiv3\bmod4),\\[2mm]
\dfrac12&(\ell\equiv1\bmod4),\end{cases}\]
where $\chi_\ell$ is the Legendre symbol and $h(-\ell)$ the class number of $\Q(\sqrt{-\ell})$. For $\ell=7$ one has
$h(-7)=1$, so the average is $1/3$ if $n_u$ is a quadratic residue and $2/3$ otherwise.
\end{proposition}

\section{Proofs}

\subsection{The lattice lemmas}
\begin{lemma}\label{l:orbit}\tagP\
Let $L$ be a free $\Zp$-module of finite rank, $M\subseteq L\otimes\Qp$ a $\Zp$-lattice, and
$\kappa=\max\{k:M\subseteq p^kL\}$. Then the $\Zp$-submodule spanned by $\bigcup_{g\in\GL(L)}g(M)$ is $p^{\kappa}L$.
\end{lemma}
\begin{proof}
Since $g(M)\subseteq g(p^\kappa L)=p^\kappa L$, the span lies in $p^\kappa L$. Conversely, by maximality of $\kappa$
there is $x\in M\setminus p^{\kappa+1}L$, and $y=p^{-\kappa}x$ lies in $L\setminus pL$. By Nakayama's lemma $y$
extends to a basis of $L$, so for every vector $e$ of a basis of $L$ there is $g\in\GL(L)$ with $g(y)=e$. Then
$p^\kappa e=g(x)\in g(M)$. The vectors $p^\kappa e$ span $p^\kappa L$.
\end{proof}

\begin{lemma}\label{l:span}\tagP\
For $U\subseteq T_c$ and a unit $b$, $U\subseteq b\,\Ocal$ if and only if $\operatorname{span}_{\Zp}U\subseteq b\,\Ocal$.
Hence $\hull(U)=\hull(\operatorname{span}_{\Zp}U)$.
\end{lemma}
\begin{proof}
$b\,\Ocal$ is a $\Zp$-submodule.
\end{proof}

\begin{lemma}\label{l:tensor}\tagP\
The $\Zp$-submodule spanned by $\Theta_{i,p,c}$ is $p^{\kappa}\bigotimes\I$, where $\kappa=\kappa_{v_j}(\theta)$.
\end{lemma}
\begin{proof}
Since $\theta$ acts on the factor of label $j$,
$\theta\cdot\bigotimes\I=\I_{v_1}\otimes\cdots\otimes\theta\I_{v_j}\otimes\cdots$. For
$\phi=\bigotimes_lg_l\in\mathrm{tensorAut}$ one has $g_l(\I_{v_l})=\I_{v_l}$, so
$\phi(\theta\cdot\bigotimes\I)=\I_{v_1}\otimes\cdots\otimes g_j(\theta\I_{v_j})\otimes\cdots$. Writing $A$ for the free
module formed by the other factors, the span of the union of these images is
$A\otimes\operatorname{span}\bigl(\bigcup_{g_j}g_j(\theta\I_{v_j})\bigr)$. By Lemma~\ref{l:orbit} with $L=\I_{v_j}$ and
$M=\theta\I_{v_j}$, this is $A\otimes p^\kappa\I_{v_j}=p^\kappa\bigotimes\I$.
\end{proof}

\subsection{Proof of Theorem~\ref{t:closed}}
By Lemmas~\ref{l:span} and~\ref{l:tensor}, $\hull(\Theta_{i,p,c})=\hull(p^\kappa\bigotimes\I)$. The hull commutes with
multiplication by units (lana, \lean{leastHull\_units\_smul}), so this is $p^\kappa\,\hull(\bigotimes\I)$. Put
$X=\hull(\bigotimes\I)$. Log-volumes of unit multiples shift additively (\lean{componentVol\_units\_smul}):
multiplication by $p^\kappa$ shifts by $-\kappa\log p$, and multiplication by $\theta$ by $-\ordp(\theta)\log p$
(\lean{componentVol\_prod\_incl\_smul\_integral}). Hence
\[\itwo=\bigl(\mulog(X)-\kappa\log p\bigr)-\bigl(\mulog(X)-\ordp(\theta)\log p\bigr)=(\ordp\theta-\kappa)\log p.\qedhere\]

\subsection{Proof of Theorem~\ref{t:bounds}}
\emph{Lower bound.} Compare the Haar measures of the two sides of $\theta\I_v\subseteq p^\kappa\I_v$ in $K_v$. With
$n=[K_v:\Qp]$, multiplication by $\theta$ multiplies measures by $p^{-n\ordp\theta}$ and multiplication by $p^\kappa$
by $p^{-n\kappa}$, so $\kappa\le\ordp\theta$. As $\kappa$ is an integer, $\kappa\le\lfloor\ordp\theta\rfloor$, and
Theorem~\ref{t:closed} gives
\[\itwo\ge(\ordp\theta-\lfloor\ordp\theta\rfloor)\log p=\fr(\ordp\theta)\log p.\]
\emph{Upper bound.} IUT~IV, Proposition~1.2(i), gives
$\theta\I_v\subseteq p^{\lfloor\ordp\theta-a_v-b_v\rfloor}\I_v$ (lana, \lean{exists\_mul\_mem\_logShell}). Hence $\kappa\ge\lfloor\ordp\theta-a_v-b_v\rfloor$,
and Theorem~\ref{t:closed} gives the upper bound; $x-\lfloor x-c\rfloor<1+c$ gives the last inequality.

\emph{Small places.} Here $a_v+b_v=0$ (IUT~IV, Proposition~1.2; lana, \lean{logLatticeA\_add\_B\_eq\_zero}), so the
two bounds coincide.\qed

\subsection{Proof of Theorem~\ref{t:pos}}
(i) $\ordp\theta=j^2\ordp(\qb_v)=j^2\operatorname{ord}_w(q_w)/(2\ell e_w)$. If this were an integer $m$, then
$j^2\operatorname{ord}_w(q_w)=2\ell e_wm$, and the prime $\ell$ would divide $j^2\operatorname{ord}_w(q_w)$. But
$1\le j\le\ell^*<\ell$, so $\ell\nmid j$, and $\ell\nmid\operatorname{ord}_w(q_w)$ (\lean{q\_order\_coprime}). Hence
$\fr(\ordp\theta)>0$, and $\itwo>0$ by Theorem~\ref{t:bounds}.

(ii) Each term of the formula for $\Delta(D)$ is the product of $w_c>0$ (\lean{weight\_pos}) and $\itwo\ge0$
(\lean{ind2Term\_nonneg}). Since $\Vbad$ is nonempty, take $u\in\Vbad$, let $p$ be its residue characteristic and $c$
a tuple with $c(j)=u$; by (i) that term is positive.\qed

\subsection{Proofs of Corollary~\ref{c:size} and Proposition~\ref{p:classno}}
\emph{Corollary~\ref{c:size}.} By Theorem~\ref{t:closed} the term depends only on $c(j)$. Since the weights over $p$
sum to $1$, $\sum_cw_c\,f(c(j))=\sum_uw_u\,f(u)$. At places that are not bad, $\qb=1$ and the term is $0$. This gives
the formula; the lower bound and the case of equality follow from Theorem~\ref{t:bounds}. The upper bound follows from
the bound $(1+a_v+b_v)\log p$ of Theorem~\ref{t:bounds} and from $(a_v+b_v)\log p\le3+\log e_v$ (lana,
\lean{logLatticeA\_add\_B\_le}).\qed

\emph{Proposition~\ref{p:classno}.} As $j$ runs through $1,\dots,\ell^*$, $j^2\bmod\ell$ runs through the set $Q$ of
quadratic residues, each once. The sum is therefore $\frac1\ell\sum_{x\in nQ}x$ (representatives in
$1,\dots,\ell-1$), and $nQ$ is $Q$ if $n$ is a residue and the set $N$ of non-residues otherwise. If
$\ell\equiv1\bmod4$, then $-1\in Q$ and $\sum_Qx=\ell(\ell-1)/4$. If $\ell\equiv3\bmod4$ and $\ell>3$, Dirichlet's
class number formula $h(-\ell)=\frac1\ell\bigl(\sum_Nx-\sum_Qx\bigr)$ and $\sum_Q+\sum_N=\ell(\ell-1)/2$ give
$\sum_Qx=\ell(\ell-1)/4-\ell h(-\ell)/2$. Divide by $\ell^*$.\qed

\section{Examples}\label{s:examples}
\tagP\ We computed the six explicit initial $\Theta$-data with $\ell=7$ of the author's catalogue~\cite{Cat}
($\ell^*=3$, $s=\ell(\ell+1)/12=14/3$); its seventh datum, $D_{2357}^{(\ell)}$, has $\ell\ge89$ and is not treated
here. The catalogue takes as $\Vbad$ the places of multiplicative reduction outside $2\ell$. We take this set minus
the places above $3$ and $5$, since above $3$ and $5$ the ramification of $K$ depends on the construction of $F$; in
these data no other place is left out, and $\ell\nmid\operatorname{ord}q$ at every place we keep. IUT~I,
Definition~3.1(b), requires $\Vbad$ only to be a nonempty set of nonarchimedean places of odd residue characteristic
such that $X_F$ has bad (multiplicative) reduction at the places of $F$ over it, and notes that $X_F$ may have bad
reduction over $\mathbb{V}_{\mathrm{mod}}^{\mathrm{good}}$ as well; Definition~3.1(c) requires $\ell$ to be prime to
the residue characteristics of $\Vbad$ and to the $q$-orders there. So this choice is allowed~\cite{IUT1}. The lower bound is the sum of
fractional parts in Theorem~\ref{t:bounds} and depends only on the data. The upper bound uses the constants $a,b$ of
IUT~IV, Proposition~1.2. For the ramification index of $K_v/\Qp$ we take
$e(F_{\mathrm{mod},u}/\Qp)\cdot15\ell/\gcd(15\ell,\operatorname{ord}_uq)$; above the special prime this is the formula
of Appendix~A of the catalogue, at the other places it is an assumption. The assumption affects only which places are
small and the upper bound.

\begin{center}
\begin{tabular}{lcrrrr}
\toprule
datum & bad places & lower bound & upper bound & $(s-1)|\log q|$ & ratio\\
 & (small, not small) & of $\Delta(D)$ & of $\Delta(D)$ & & \\
\midrule
D211 & 4, 1 & 25.894 & 27.933 & 79.937 & 0.324--0.349\\
D109 & 6, 0 & 22.707 & 22.707 (exact) & 70.085 & 0.324\\
D2 & 3, 2 & 28.516 & 30.914 & 62.881 & 0.453--0.492\\
D3 & 14, 5 & 93.174 & 102.691 & 295.492 & 0.315--0.348\\
D18 & 14, 3 & 53.621 & 58.621 & 161.059 & 0.333--0.364\\
D15 & 13, 2 & 45.765 & 49.165 & 135.160 & 0.339--0.364\\
\bottomrule
\end{tabular}
\end{center}

\begin{example}[D109]\tagP\
There are six bad places, with residue characteristics $109$, $21577$, $6065177$, $534775081$,
$1174770007439769673$ and $16450055877403724897$. All weights are $1/2$, and $r_u$ is $29/7$ at $109$ and $1/7$ at the
others. Even at the smallest, $p=109$, one has $e_v=105\le p-2$ (by the formula of Appendix~A of the catalogue), so all
six places are small. Both $n_u=29$ and $n_u=1$ are quadratic residues modulo $7$, so by Proposition~\ref{p:classno}
the label average is $1/3$ and
\[\Delta(D_{109})=\sum_u\frac12\cdot\frac13\log p_u=\frac16\sum_u\log p_u=22.7068\ldots\]
The ratio to $(s-1)|\log q|=70.085$ is $0.324$.
\end{example}

\begin{example}[a model of a place that is not small]\tagP\
Let $K_v=\Q_3(\varpi)$ with $\varpi^4=3$ ($p=3$, $e=4$). We computed $\I_v$ and $\kappa$ exactly. For $\theta=\varpi$,
$\ordp\theta=1/4$ and $\kappa=-1$, so the term is $\tfrac54\log3$, which exceeds $\log p$. For $\theta=\varpi^4$,
$\kappa=1$ and the term is $0$; for $\theta=(1+\varpi)\varpi^4$, $\kappa=0$ and the term is $\log3$: the value
depends on the unit part. All three lie inside the upper bound $(\ordp\theta-\lfloor\ordp\theta-a-b\rfloor)\log p$ of
Theorem~\ref{t:bounds}, here with $a+b=7/4$. This model is for illustration only and is not a bad place of actual data
(at a bad place $\ordp\theta\notin\Z$).
\end{example}

\tagH\ In these data the $q$-order is $2$ at most bad places, which is small, so $\Delta(D)$ is between three and five
tenths of the main term $(s-1)|\log q|$. That $\Delta(D)$ is small holds only in the direction of growing $q$-orders;
its actual size is set by the values $\log p$ at the bad places.

\section{Discussion}

\subsection{How the two walls relate}
\tagP\ With lana's identity $\Delta(D)=R+s|\log q|-V$, where $V=\mulog(\hull(VC(I)))$ (\lean{gap\_eq\_rhs\_add}), the
two walls read
\[\text{(SS)}\quad(s-1)|\log q|\le V,\qquad\qquad\text{(Var)}\quad(s-1)|\log q|\le V+\Delta(D).\]
\tagH\ $V$ is bounded above in terms of ramification and does not involve the $q$-orders (lana,
\lean{Iut/WallEquiv/VcI.lean}); by Corollary~\ref{c:size}, $\Delta(D)$ is a quantity of the same kind. So the two
walls are height inequalities of the same shape that differ only in a term of conductor type on the right-hand side.
The difference is positive (Theorem~\ref{t:pos}).

\subsection{Readings of the (Ind2) group}\label{s:groups}
\begin{itemize}
\item \tagP\ IUT~IV, Proposition~1.2, considers arbitrary $\Qp$-linear automorphisms $\phi$ preserving
$\bigotimes\I$, a larger group than $\mathrm{tensorAut}$. The value does not change for that group: since $\theta$
acts on one factor only, $\theta\cdot\bigotimes\I\subseteq p^k\bigotimes\I$ if and only if
$\theta\I_{v_j}\subseteq p^k\I_{v_j}$, so Lemma~\ref{l:orbit} applied to $L=\bigotimes\I$ yields the same $\kappa$.
\item \tagP\ For a smaller group $G\subseteq\mathrm{tensorAut}$ the union only shrinks, so the term lies between $0$
and the value found here. The image of the group $\mathrm{Ism}$ of the original papers is such a case (lana's
documentation).
\item \tagP\ At a small place, if $G$ consists only of multiplications by $\Zp^{\times}$ and field automorphisms, the
term is $0$, because $\sigma(\theta)\I_v=\theta\I_v$. If moreover every bad place is small, then $\Delta(D)=0$ and
(SS) and (Var) coincide.
\item \tagH\ Where the value for the group of the original papers lies in this range depends on how far that group
moves $\theta\I_v$. The proof of Lemma~\ref{l:orbit} uses elements that send a primitive element of $\theta\I_v$ to a
basis vector. Whether the image of $\mathrm{Ism}$ contains such elements is not decided here.
\item \tagP\ The author's fourth report~\cite[\S6, Proposition~9]{Fourth} applies \cite[Thm.~F, Thm.~7.4]{MT} to a
finite extension $K_v$ of $\Qp$: there $\mathrm{Ism}(G_{K_v})=\Zp^{\times}$, so (Ind2) acts on each factor by a unit
scalar. If the group of the original papers is identified with $\mathrm{Ism}$ in this way, it fixes every
$\theta\I_v$, since $\I_v$ is a $\Zp$-module; the union $\Theta_{i,p,c}$ is then $\theta\cdot\bigotimes\I$, and its
(Ind2)-term is $0$, because $\hull(\theta X)=\theta\,\hull(X)$ (proof of Theorem~\ref{t:closed}). For that group
$\Delta(D)=0$, and (SS) and (Var) coincide.
\end{itemize}

\subsection{Comparison with the estimate of IUT~IV}
\tagS\ IUT~IV, Proposition~1.2(iii) states: if $p>2$ and $e_i\le p-2$ for all $i\in I$, then
$\phi(p^\lambda\cdot(R_I)^{\sim})\subseteq p^{\lambda-\mathfrak{d}_I-1}\cdot(R_I)^{\sim}$ for every
$\lambda\in\frac1{e_i}\Z$, $i\in I$~\cite{IUT4}.

\tagH\ Besides the different term $\mathfrak{d}_I$, this estimate allows $1$ for rounding. The term $\mathfrak{d}_I$
comes from comparing $\bigotimes_iR_i$ with its normalization $(R_I)^{\sim}$ (the proof of (ii) uses Proposition~1.1
for this). In lana's $\Delta(D)$ the hull is taken on both sides, so this comparison, the shell defect, cancels (the
docstring of \lean{gap\_eq}; for the shell defect see \lean{Iut/WallEquiv/VcI.md}), and only the rounding part
remains. Theorem~\ref{t:bounds} shows that under lana's group this rounding loss actually occurs, and at small places
its size is exactly $\fr(\ordp\theta)\in(0,1)$ in units of $\log p$. Averaged over the labels it is about $1/2$ by
Proposition~\ref{p:classno} ($1/3$ or $2/3$ for $\ell=7$). Compared with the $1$ that the estimate of IV allows for
rounding, the actual loss is about half of it on average over the labels.

\tagP\ The unused part of this allowance, $1-\fr(\ordp\theta)=1+\lfloor\ordp\theta\rfloor-\ordp\theta$ at small
places, is the floor slack that the author's third report computes on explicit data~\cite[\S6]{Third}.

\tagX\ Zeng~\cite{Zeng} (first version 29 October 2026) computes the index by which the hull enlarges the integral
structure of a tensor product of rings of integers. \tagP\ In $\Delta(D)$ both terms are log-volumes of unit multiples
of the same hull $X=\hull(\bigotimes\I)$ (proof of Theorem~\ref{t:closed}), so the enlargement of $\bigotimes\I$ to its
hull, of the kind Zeng measures, cancels in $\Delta(D)$.

\section{Reproducibility}
The folder \texttt{ind2-gap/} of the repository
\url{https://github.com/turahigashi/iut-explicit-initial-theta} (tag \texttt{v0.24.0}), also deposited with this note,
contains the following scripts and their outputs. They use the program \texttt{iut\_compute.py} of the author's
walkthrough~\cite{Walk}, which is in the same repository.
\begin{itemize}
\item \texttt{scripts/check\_delta.py}: the table of Section~\ref{s:examples} (lower bounds by fractional parts, upper
bounds by the constants $a,b$), and Proposition~\ref{p:classno} checked in exact rational arithmetic for twelve primes
$\ell\equiv3\bmod4$ ($7\le\ell\le83$) and six primes $\ell\equiv1\bmod4$ ($5\le\ell\le41$), for all $n$.
\item \texttt{scripts/delta2.py}: the same lower bounds, with upper bounds from the inner and outer radii of the
log-shells.
\item \texttt{scripts/shell\_lattice.py}, \texttt{scripts/gl\_check.py}: the log-shell $\I_v=(2p)^{-1}\log(1+\mathfrak m)$
of $\Qp(p^{1/e})$ in exact rational arithmetic, and the span of $\theta\I_v$ under random elements of $\GL(\I_v)$
compared with $p^\kappa\I_v$ (Lemma~\ref{l:orbit}, Example~5.2) for $(p,e)=(11,3),(3,4),(5,6)$.
\end{itemize}

\section*{Acknowledgements and use of AI}
The question studied in this note---to determine the value of the (Ind2)-terms that lana-agents/iut leaves
open---arose in the author's discussions with generative AI tools; the author judged it important and decided to make
it the subject of this note. Its ultimate aim, to understand exactly the algorithm by which IUT computes the height
inequality, is the author's research direction. Generative AI tools (large language models) were used extensively,
under the author's direction, in finding the proofs, in the computations and in the writing. The author takes full
responsibility for the contents of this note.

\begin{thebibliography}{IUT4}\raggedright
\bibitem[Cat]{Cat}
Toshihisa Urahigashi, \emph{Explicit initial $\Theta$-data and the two readings of Step (v) of IUT IV, Theorem 1.10:
a catalogue}, version 1.0.0, Zenodo, 2026, \href{https://doi.org/10.5281/zenodo.23870448}{doi:10.5281/zenodo.23870448}.

\bibitem[DH1]{DH1}
Taylor Dupuy and Anton Hilado, \emph{The statement of {M}ochizuki's {C}orollary 3.12, initial theta data, and the
first two indeterminacies}, arXiv:2004.13228v1 [math.NT], 28 April 2020; the section cited (\S4.8) is in version~1
(version~2 of 2025 is retitled and has no \S4).

\bibitem[DH2]{DH2}
Taylor Dupuy and Anton Hilado, \emph{Probabilistic {S}zpiro, baby {S}zpiro, and explicit {S}zpiro from {M}ochizuki's
{C}orollary 3.12}, arXiv:2004.13108 [math.NT], 2020.

\bibitem[Fourth]{Fourth}
Toshihisa Urahigashi, \emph{Split bad primes and the tuple structure of inter-universal {T}eichm\"uller theory: what
the source states, and what follows (a fourth report)}, version 1.0.4, Zenodo, 2026,
\href{https://doi.org/10.5281/zenodo.22648944}{doi:10.5281/zenodo.22648944}.

\bibitem[IUT1]{IUT1}
Shinichi Mochizuki, \emph{Inter-universal {T}eichm\"uller theory {I}: construction of {H}odge theaters}, Publ. Res.
Inst. Math. Sci. \textbf{57} (2021), 3--207.

\bibitem[IUT3]{IUT3}
Shinichi Mochizuki, \emph{Inter-universal {T}eichm\"uller theory {III}: canonical splittings of the log-theta-lattice},
Publ. Res. Inst. Math. Sci. \textbf{57} (2021), 403--626.

\bibitem[IUT4]{IUT4}
Shinichi Mochizuki, \emph{Inter-universal {T}eichm\"uller theory {IV}: log-volume computations and set-theoretic
foundations}, Publ. Res. Inst. Math. Sci. \textbf{57} (2021), 627--723.

\bibitem[lana]{lana}
lana-agents, \emph{iut}: a Lean~4 formalization of inter-universal {T}eichm\"uller theory, GitHub repository,
\url{https://github.com/lana-agents/iut}, commit \texttt{553076a} (3 November 2026).

\bibitem[LANA]{LANA}
Project LANA, \emph{Project LANA interim report on IUT theory}, 20 July 2026,
\url{https://ncatlab.org/nlab/files/LANAProject-Report-July2026.pdf}.

\bibitem[MT]{MT}
Arata Minamide and Shota Tsujimura, \emph{Anabelian geometry for {H}enselian discrete valuation fields with
quasi-finite residues}, RIMS Preprint 1973, June 2023.

\bibitem[SS]{SS}
Peter Scholze and Jakob Stix, \emph{Why abc is still a conjecture}, 2018.

\bibitem[Third]{Third}
Toshihisa Urahigashi, \emph{The cost of each distinction in the local computations of inter-universal
{T}eichm\"uller theory~{IV} (a third report)}, version 1.0.2, Zenodo, 2026,
\href{https://doi.org/10.5281/zenodo.22646136}{doi:10.5281/zenodo.22646136}.

\bibitem[Walk]{Walk}
Toshihisa Urahigashi, \emph{A programmer's walkthrough of inter-universal {T}eichm\"uller theory}, version 1.0.0,
Zenodo, 2026, \href{https://doi.org/10.5281/zenodo.23880260}{doi:10.5281/zenodo.23880260}.

\bibitem[Zeng]{Zeng}
Chunping Zeng, \emph{An integral-structure inflation in the hull operation of {IUT} {C}orollary 3.12}, Zenodo, 2026,
first version 29 October 2026, \href{https://doi.org/10.5281/zenodo.23754276}{doi:10.5281/zenodo.23754276}.
\end{thebibliography}

\end{document}
